% resumo em português
\setlength{\absparsep}{18pt} % ajusta o espaçamento dos parágrafos do resumo

\begin{resumo}
 Este trabalho apresenta uma análise comparativa do desempenho de redes Wi-Fi corporativas (Ruckus e UniFi) utilizando telemetria física e lógica. Com base em um dataset de 18 dias coletado em um campus universitário de alta densidade, contendo mais de 350 mil registros, o estudo propõe um pipeline estruturado de limpeza, testes estatísticos de hipóteses e classificação heurística e não-supervisionada de gargalos. Testes de Mann-Whitney revelaram disparidades significativas ($p < 0,001$) entre as bandas de 2,4~GHz e 5~GHz, e taxas de retransmissão discrepantes entre os fabricantes (UniFi com 22,1\% vs. Ruckus com 5,7\% em média). A classificação de gargalos indicou que o throughput dos clientes foi severamente atenuado por sinal fraco (degradação de até 89,3\% na vazão) e instabilidade de roaming, que reduziu a vazão no Ruckus em 38,8\%. O algoritmo K-Means ($K=4$) foi validado pelas métricas de Silhueta (0,3383) e Davies-Bouldin (0,7953) para segmentação dos Access Points em perfis operacionais. A pesquisa contribui com uma metodologia reprodutível para o diagnóstico automatizado de gargalos em redes sem fio empresariais.

 \textbf{Palavras-chave}: Telemetria de Rede. Wi-Fi Corporativo. K-Means. Análise Estatística. Diagnóstico de Gargalos.
\end{resumo}
